\Needspace{1.25in}
\manualtablecaption{Atomic Memory-Order Selectors}
\begingroup\footnotesize
\setlength{\aboverulesep}{0pt}
\setlength{\belowrulesep}{0pt}
\setlength{\extrarowheight}{1.2pt}
\begin{center}
\begin{tabularx}{\linewidth}{@{}>{\raggedright\arraybackslash}p{0.95in}>{\raggedright\arraybackslash}p{0.45in}>{\raggedright\arraybackslash}X@{}}
\toprule
\rowcolor{ManualHeaderFill}
\textbf{Selector} & \textbf{Code} & \textbf{Architectural Ordering Effect}\\
\midrule
\texttt{RELAXED} & \texttt{0} & Atomicity only. No ordering is imposed on unrelated memory operations.\\
\texttt{ACQUIRE} & \texttt{1} & Later memory operations by the same logical processor are ordered after the atomic operation.\\
\texttt{RELEASE} & \texttt{2} & Earlier memory operations by the same logical processor are ordered before the atomic operation.\\
\texttt{ACQREL} & \texttt{3} & Applies both acquire and release ordering.\\
\texttt{SEQCST} & \texttt{4} & Applies acquire-release ordering and participates in the single global sequentially consistent order.\\
\bottomrule
\end{tabularx}
\end{center}
\endgroup
